\subsection{Local rational root bounds near the indexed resolution edge}
\label{ssec:local-rational-root-bound}

The global Fourier derivative bound in
Section~\ref{ssec:quantitative-mode-transfer} is deliberately conservative.
It can fail to establish a useful neighborhood when the transform and its
derivatives are small together. The following local bound retains that
cancellation in the Taylor coefficients. It is elementary rational-function
analysis, not a new asymptotic formula for the arithmetic ground state.

For an even finite state, set $F_u(t)=p(t)R_u(t)$, with
$p(t)=2L^{-1/2}\sin(tL/2)$. Away from the Fourier lattice,
\[
 R_u(t)=\frac{u_0}{t}+\frac1{\sqrt2}\sum_{j=1}^N u_j
 \left\{\frac1{t-\alpha_j}+\frac1{t+\alpha_j}\right\},
 \qquad \alpha_j=2\pi j/L.
\]
Here $u_j$ are orthonormal even-sector coefficients. Let
$R(t)=\sum_i w_i/(t-p_i)$, fix a real center $x$ off its poles, and put
\[
 c_m=(-1)^m\sum_i\frac{w_i}{(x-p_i)^{m+1}},\qquad
 d_i=|x-p_i|,\qquad q_i=r/d_i<1.
\]
For an integer $M\ge1$ define
\begin{align*}
 V_M(r)&=\sum_{m=2}^M m|c_m|r^{m-1}+T_M(r),\\
 T_M(r)&=\sum_i\frac{|w_i|}{d_i^2}\,
 \frac{q_i^M\{(M+1)-Mq_i\}}{(1-q_i)^2}.
\end{align*}

\begin{proposition}[Local derivative margin and displacement enclosure]
\label{prop:local-rational-root-bound}
For $|h|\le r<\min_i d_i$,
$|R'(x+h)-c_1|\le V_M(r)$. If
\[
 a:=|c_1|>V_M(r),\qquad |c_0|<\{a-V_M(r)\}r,
\]
then $R$ has exactly one real root $x+h_*$ in $(x-r,x+r)$, and
\[
 \left|h_*+\frac{c_0}{c_1}\right|\le\frac{V_M(r)}a\,r.
\]
If this interval also excludes every zero of $\sin(tL/2)$, the same
existence and uniqueness hold for the finite transform in that interval.
\end{proposition}
\begin{proof}
Expand each simple fraction geometrically. The differentiated tail is
bounded by
\[
 \frac{|w_i|}{d_i^2}\sum_{m=M+1}^\infty m q_i^{m-1}
 =\frac{|w_i|}{d_i^2}
 \frac{q_i^M\{(M+1)-Mq_i\}}{(1-q_i)^2}.
\]
Keep the cancellations in the finitely many $c_m$ before taking absolute
values. The derivative has constant sign and magnitude at least
$a-V_M(r)$ throughout the interval. The strict residual inequality
gives opposite endpoint signs and therefore one root. Integrating
$R'(x+h)-c_1$ from zero to that root gives the displacement bound.
The last assertion follows because the sine factor never vanishes there.
\end{proof}

All quantities admit outward-rounded interval evaluation. An unresolved
derivative, a radius reaching a pole or a carrier zero, or an unsuccessful
Taylor-order limit is an unresolved certificate, not a claimed absent
root. Arithmetic precision used to verify this finite rational function
is distinct from accuracy of its source eigenvector.

\paragraph{A stronger finite check on the retained mode pair.}
The unchanged $c=1000$, $N=800\to890$ pair from
Section~\ref{ssec:quantitative-mode-transfer} was evaluated at all thirteen
previously frozen indices. The parent and child vectors are the exact
decimal vectors in the retained files; no eigensolve or matrix assembly
was repeated. Interval evaluation uses 2100 decimal digits, and an
independent direct-power evaluation at 2600 digits rechecks every certified
derivative margin and both endpoint signs. Nine hundred separate exact
rational fixture checks verify the derivative-tail implementation.

All thirteen parent neighborhoods and nine child neighborhoods certify.
The nine originally eligible pairs all retain the one-percent pass for
the frozen full two-term Newton prediction. The largest certified
relative comparison bound is below $7.36\cdot10^{-10}$, including the
deliberate enclosure padding; it is not an estimate of the actual model
error. The old global energy/Fourier-radius calculation established
neighborhoods at only four of these indices. The new child orders are
\begin{center}\small
\begin{tabular}{@{}rrrrrrrrrr@{}}
\toprule
$k$ & 1 & 20 & 100 & 300 & 400 & 450 & 475 & 500 & 525\\
$M$ & 1 & 1 & 1 & 1 & 1 & 3 & 5 & 9 & 21\\
\bottomrule
\end{tabular}
\end{center}
The four originally ineligible child cases remain unresolved: $k=550$
reaches the fixed order limit 96, while the proposed neighborhoods for
$k=575,600,610$ meet a Fourier carrier zero. They are retained rather
than removed or relabeled as converged. These are local certificates for
serialized states, not analytic-ground source certificates, full-spectrum
ordinal counts, or thirteen independent matrix experiments.

\paragraph{Three different notions of sufficient resolution.}
The condition $\alpha_{N+1}>\nu$ in the earlier adjoint-tail bound is a
condition for that particular estimate, not a necessary condition for an
accurate movable root. In this retained example, $N=800,k=500$ has about
73.905 relative digits while $\nu L/(2\pi)\simeq891.819>N+1$.
A hard rule discarding every target beyond the retained Fourier band
would therefore discard a successful reproduction.

There is nevertheless a degree obstruction: after multiplication by
$t\prod_{j=1}^N(t^2-\alpha_j^2)$, the secular numerator is a polynomial
of degree at most $N$ in $t^2$. A nonzero even state therefore has at
most $N$ distinct positive nonlattice secular roots. Resolving an entire
prefix of $K$ such roots requires $N\ge K$; it does not guarantee the
required accuracy, nor identify a reference label with a spectral ordinal.

Finally, writing $\rho=\nu/\alpha_{N+1}<1$, the earlier multiplier is
\[
 \frac{\nu}{\alpha_{N+1}^2-\nu^2}
 =\frac1{\alpha_{N+1}}\frac{\rho}{1-\rho^2}.
\]
Its fixed-$k$ $N^{-2}$ behavior becomes only $N^{-1}$ when $\rho$ stays
strictly between zero and one, and loses that suppression as
$\alpha_{N+1}-\nu$ remains bounded. These statements concern this explicit
factor at fixed cutoff, with the other theorem constants controlled;
they are not lower bounds for the actual error. An accuracy-dependent
minimum-$N$ law must distinguish the degree floor, an estimate's
admissibility conditions, and the actual indexed accuracy threshold.

\paragraph{What this removes, and what remains.}
This local certificate closes the root-nonlinearity qualification for
the nine retained eligible comparisons. It does not determine the
arithmetic mode innovation, cutoff displacement, or source uncertainty
from $c,N,P,k$. Those remain the work needed for the parameter law.
